\documentclass[10pt,a4paper]{amsart}
\usepackage[margin=19mm]{geometry}
\usepackage{amsmath,amssymb,mathtools}
\usepackage[scr=rsfs,cal=euler]{mathalfa}
\usepackage{xfrac}
\usepackage[unicode,hidelinks,bookmarks=false]{hyperref}
\newcommand{\ie}{i.e.\ }
\newcommand{\cf}{cf.\ }
\newcommand{\resp}{respectively\ }
\newcommand{\Subsection}{\subsection}
\providecommand{\nobreakdash}{\nobreak\hbox{-}}
\setlength{\emergencystretch}{2em}
\multlinegap0pt
\usepackage{bm}
\usepackage{bbm}
\usepackage{dsfont}
\usepackage{xspace}
\usepackage{cancel}
\usepackage{mathtools}
\usepackage[shortlabels]{enumitem}
\usepackage{amsthm}
\makeatletter
\@ifundefined{theorem}{\newtheorem{theorem}{Theorem}}{}
\@ifundefined{fact}{\newtheorem{fact}{Fact}}{}
\@ifundefined{proposition}{\newtheorem{proposition}[theorem]{Proposition}}{}
\makeatother
\makeatletter
\expandafter\ifx\csname claimproof\endcsname\relax
\newenvironment{claimproof}[1][Proof of claim]{\begin{proof}[#1]\renewcommand{\qedsymbol}{$\scriptstyle\square$}}{\end{proof}}
\fi
\def\tc{\mathop{\text{\rm tc}}\nolimits}
\makeatother
\title[Monochromatic components in dense 2-edge-coloured balanced b]{Monochromatic components in dense 2-edge-coloured balanced bipartite graphs}
\author{C\'esar Bispo, George Kontogeorgiou, Marcelo Lage, Guilherme O. Mota, Bruno Skarmeta}
\date{Source: arXiv:2608.17300v1 (CC BY 4.0)}
\begin{document}
\maketitle

\noindent\emph{Source and licence.} Monochromatic components in dense 2-edge-coloured balanced bipartite graphs. Version arXiv:2608.17300v1 (CC BY 4.0), distributed under the Creative Commons Attribution 4.0 International licence, as recorded in the arXiv licence metadata for this exact version and shown on the abstract page https://arxiv.org/abs/2608.17300v1. Full rights notice: licenses/arxiv-2608.17300-CC-BY-4.0.txt.\par\smallskip

\noindent\textit{Editorial scope.} Counted unit: the Proposition whose source label is \texttt{prop:7} in the article (source0.tex lines 699--711, source label \texttt{prop:7}), delivered together with the article's own complete original proof of it (source0.tex lines 713--829, one \texttt{proof} environment of 117 lines). No mathematics is written, repaired or re-derived: statement and proof are the article's own text line for line. Required context retained uncounted: fact:simple (lines 376--385); prop:5 (lines 479--486). Every definition, display and label of those lines is retained unchanged. Omitted from this package and not redistributed: 1--375, 386--478, 487--698, 712--712, 830--1623. Nothing inside a retained excerpt is omitted or altered: every retained body is delivered exactly as the article's lines are written, so no omission receipt of any kind accompanies this package. Citations: the retained excerpts carry no citation command at all, so no bibliography entry of the article is transcribed and no reference is redistributed. Reference adaptation: none was needed and none was made; every reference inside the delivered unit resolves inside it, so no unresolved reference or citation is produced. Printed slips kept literal: no printed word, symbol, hypothesis or number of the counted statement or of its proof was altered. Layout and class: the article's own class is \texttt{amsart} with packages including lmodern, fullpage, setspace, datetime, verbatim, url, dsfont, xcolor, tikz, amsmath, cite, euflag, graphicx, caption; the wrapper is the lane's pinned \texttt{amsart} editorial wrapper at 10pt on A4, which restates the theorem-like environments the retained text needs and adds the article's own macro definitions that the retained text needs (\texttt{\textbackslash tc}), with extra wrapper packages (bm, bbm, dsfont, xspace, cancel, mathtools, enumitem). The delivered numbering is the unit's own. Assets: no retained span contains a figure environment, a graphics-inclusion command, a PDF-page inclusion, a table or a TikZ picture; no image file is copied into this package, the package declares no asset dependency and no image file is redistributed with it. Licence: version arXiv:2608.17300v1 of this article is distributed under the Creative Commons Attribution 4.0 International licence, as recorded in the arXiv licence metadata for this exact version and shown on the abstract page https://arxiv.org/abs/2608.17300v1. A full rights notice is delivered at licenses/arxiv-2608.17300-CC-BY-4.0.txt. Characters: every retained excerpt is pure ASCII, so the delivered engine, XeLaTeX with its default Latin Modern text font, reports no missing character.
\begin{fact}\label{fact:simple}
For any positive integer $n$, we have 
\[
\left\lceil\frac{\delta^*(n)}{2}\right\rceil+\delta^*(n)\ge n,
\qquad
3\left\lceil\frac{\delta^*(n)}{2}\right\rceil\ge n,
\qquad
3\left(\left\lfloor\frac{\delta^*(n)}{2}\right\rfloor+1\right)>n.
\]
\end{fact}

\begin{proposition}\label{prop:5}
%
Let $n$ be a positive integer, and let $G=(X,Y;E)$ be a bipartite graph with
$|X|=|Y|=n$ such that $\delta(G)\ge \lfloor n/2\rfloor+1$. Let $\varphi$ be a
$2$-edge-colouring of $E(G)$. If at most two monochromatic components cover $X$ or
cover $Y$, then $\tc_2(G,\varphi)\le 3$.

\end{proposition}

\begin{proposition}\label{prop:7}
%
Let $n$ be a positive integer, and let $G=(X,Y;E)$ be a bipartite graph with
$|X|=|Y|=n$ such that $\delta(G)\ge\delta^*(n)$. Let $\varphi$ be a
$2$-edge-colouring of $E(G)$. Then the following statements hold.
\begin{enumerate}[(i)]
\item If there is a monochromatic component that covers at least $\delta^*(n)$
vertices on one side of the bipartition of $G$, then $\tc_2(G,\varphi) \le 3$;
\label{prop:7-item1}
\item If there are at most three components of the same colour that cover one
side of the bipartition of $G$, then $\tc_2(G,\varphi) \le 3$.\label{prop:7-item2}
\end{enumerate}
\end{proposition}

\begin{proof}
We start by proving item~\ref{prop:7-item1}. By symmetry, let $R_1$ be a red component
covering at least $\delta^*(n)$ vertices of $Y$. For $n\le 2$, the result is
immediate, so assume that $n\ge 3$. Since $\delta^*(n)>n/2$,
Proposition~\ref{prop:5} applies whenever one side is covered by at most two
monochromatic components.

First suppose that there exists a red-dominant vertex $u \in X \setminus
R_1$ and let $R_2$ be the red component containing $u$. Then, $|Y\cap R_2| \ge
\lceil\delta^*(n)/2\rceil$. Therefore, from Fact~\ref{fact:simple} we have
$|Y\cap(R_1\cup R_2)|\ge\delta^*(n)+\lceil\delta^*(n)/2\rceil\ge n$.
Thus two red components cover~$Y$, and Proposition~\ref{prop:5} implies
$\tc_2(G,\varphi) \le 3$.

Now suppose there is no red-dominant vertex in $A:=X\setminus R_1$. Then
every vertex $x\in A$ satisfies $d_B(x)>d_R(x)$ and since $d(x)\ge \delta^*(n)$,
we have $d_B(x)\ge \left\lfloor \delta^*(n)/2\right\rfloor+1$. Thus, every
blue component that intersects $A$ contains at least
$\left\lfloor\delta^*(n)/2\right\rfloor+1$ vertices of $Y$. Consequently, at
most two blue components intersect $A$: if three did, their pairwise disjoint
intersections with $Y$ would together contain more than $n$ vertices, which is a contradiction by Fact~\ref{fact:simple}.

If at most one blue component meets $A$, then $X$ is covered by $R_1$ and at
most one blue component and Proposition~\ref{prop:5} implies
$\tc_2(G,\varphi)\le 3$. Thus, we may assume that there are exactly two blue
components $B_u$ and $B_v$ covering $A$.

If every pair of vertices $x\in B_u\cap A$ and $z\in B_v\cap A$ has a common red
neighbour, then all vertices of $A$ lie in a single red component. Hence this
red component together with $R_1$ covers $X$, and Proposition~\ref{prop:5}
implies $\tc_2(G,\varphi)\le 3$. Thus, we may choose vertices $u\in B_u\cap A$ and $v\in
B_v\cap A$ with no common red neighbour. Since $B_u$ and $B_v$ are distinct blue
components, $u$ and $v$ have no common blue neighbour either. Therefore every
common neighbour of $u$ and $v$ is joined to them by edges of distinct colours.

Moreover, $u$ and $v$ have no common neighbour in $Y\cap R_1$: indeed, every
edge from $u$ or $v$ to a vertex of $Y\cap R_1$ is blue, because
$u,v\notin R_1$, and $u$ and $v$ have no common blue neighbour. Hence all
common neighbours of $u$ and $v$ lie in $Y\setminus R_1$.

Let $S$ denote the set of common neighbours of $u$ and $v$ in $Y \setminus R_1$
and $T$ the set of vertices of $Y \setminus R_1$ adjacent to at most one of $u$
and $v$. The next step is to show that $|T| \le 1$. For that, first observe
that we have
\begin{align*}
|S| &\ge d(u) + d(v) - n \ge 2\delta^*(n) - n,\\
|S| &= n - |Y \cap R_1| - |T| \le n - \delta^*(n) - |T|.
\end{align*}
This gives the desired bound on $|T|$, which is 
\begin{equation*}
|T| \le 2n - 3\delta^*(n) \le 1.
\end{equation*}
Denote by $S_u$ the blue neighbours of $u$ in $S$ and $S_v$ the blue neighbours
of $v$ in $S$, so that $S=S_u\mathbin{\dot\cup}S_v$.

If $T=\emptyset$, then the components $R_1$, $B_u$ and $B_v$ cover $Y$. They
also cover $X$, because $B_u$ and $B_v$ cover $X\setminus R_1$. Hence
$\tc_2(G,\varphi)\le 3$. Thus, we may assume that $T=\{t\}$.

We first show that $t$ is adjacent to at least one of $u$ and $v$. Suppose not.
Since $u$ and $v$ have no common blue neighbour, they have disjoint blue
neighbourhoods in $R_1\cap Y$. Hence
\[
    2\delta^*(n)\le d(u)+d(v)\le 2|S|+|R_1\cap Y|
        =2(n-|R_1\cap Y|-1)+|R_1\cap Y|
        \le 2n-\delta^*(n)-2,
\]
which contradicts $3\delta^*(n)>2n-2$. Therefore, without loss of generality,
suppose the sole vertex $t\in T$ is adjacent to $u$ and not to $v$.

If the edge $ut$ is blue, then $B_u$ covers $T$ and $S_u$. Since $B_v$ covers
$S_v$, the components $R_1$, $B_u$ and $B_v$ cover $Y$. They also cover $X$,
because $B_u$ and $B_v$ cover $X\setminus R_1$. Hence $\tc_2(G,\varphi)\le 3$.

It remains to consider the case where $ut$ is red. Let $R_u$ and $R_v$ be the
red components containing $u$ and $v$, respectively. Then $R_u$ covers
$T\cup S_v$, while $R_v$ covers $S_u$. Therefore the three red components
$R_1$, $R_u$ and $R_v$ cover $Y$.

If these three red components also cover $X$, then we are done. Thus, suppose
that there exists
\[
    x\in X\setminus (R_1\cup R_u\cup R_v).
\]
Since $R_1$, $R_u$ and $R_v$ cover $Y$, the vertex $x$ sends no red edge to
$Y$. Hence all edges incident to $x$ are blue. Suppose, without loss of
generality, that $x\in B_u$. Since $B_u$ and $B_v$ are distinct blue components,
$x$ has no blue neighbour in $B_v\cap Y$, and therefore no neighbour at all in
$B_v\cap Y$. Consequently,
\[
    \delta^*(n) \le d(x) \le n-|B_v\cap Y|.
\]
On the other hand, $v$ is blue-dominant, and so
\[
    |B_v\cap Y| \ge d_B(v) \ge \left\lceil \frac{\delta^*(n)}{2}\right\rceil.
\]
Thus,
\[
    \delta^*(n) \le d(x) \le n-|B_v\cap Y|
        \le n-\left\lceil \frac{\delta^*(n)}{2}\right\rceil
        \le \delta^*(n),
\]
where the last inequality follows from Fact~\ref{fact:simple}. Hence equality
holds throughout. In particular, $x$ is adjacent to every vertex of $Y\setminus
B_v$, and all these edges are blue. Therefore $Y\setminus B_v\subseteq B_u$, and
so $B_u\cup B_v$ covers $Y$. Since $B_u$ and $B_v$ cover $X\setminus R_1$, the
three components $R_1$, $B_u$ and $B_v$ cover $V(G)$. Hence $\tc_2(G,\varphi)\le
3$, which proves item~\ref{prop:7-item1}.

Item~\ref{prop:7-item2} follows from item~\ref{prop:7-item1}. Indeed, suppose
without loss of generality that at most three red components cover
$X$. If they also cover $Y$, then $\tc_2(G,\varphi)\le 3$. Otherwise, there is a vertex
$y\in Y$ outside all of them. Every edge from $y$ to $X$ is then blue, so the
blue component containing $y$ covers $N(y)$ and hence at least
$d(y)\ge \delta^*(n)$ vertices of $X$. Item~\ref{prop:7-item1} now gives
$\tc_2(G,\varphi)\le 3$, which finishes the proof.
\end{proof}

\end{document}
